% Section 5 -- Empirical Case Study (Dogfood). Pages target ~2.0.
% Defends C3 (cross-model adversarial reduces reward hacking) +
% C5 (concrete dogfood numbers from real audit logs).
% HARD RULE HR-F: every number traceable to a state file or a git command.
\section{Empirical Case Study (Dogfood)}
\label{sec:case}

We dogfooded \texttt{cursor\_manager} on three concurrent NeurIPS 2026
submissions: \emph{paper~A} (Computer-arithmetic / Physics In-Context
Learning), \emph{paper~B} (LLM Geometry under Scaling Laws), and
\emph{paper~C} (this paper). The empirical results in this section
are computed on paper~C only; paper~A and paper~B run on a separate
compute host and their audit logs are not yet available
(see~\S\ref{sec:case-setup} for the honest disclosure of which loop
ran when). All numbers in this section derive from immutable audit
artefacts (\texttt{git log} on branch
\texttt{main}, \texttt{wc -l} on the
system source, the smoke-test runner output, the lit-review sub-agent
verification report, and \texttt{state/escalations.jsonl} on the
deployment host -- all part of the open-source release,
Appendix~\ref{app:repo}); per HR-F no number here is a hand-estimate.

\subsection{Setup}
\label{sec:case-setup}

Each of the three papers is configured as one
worker entry in a single \texttt{config.toml}, sharing one running
instance of \texttt{cursor\_manager} on a single workstation. The
recommended cross-API-model adversarial configuration runs all three
roles via the codex CLI but with three \emph{different} profiles
defined in \texttt{\textasciitilde/.codex/config.toml}: the worker
uses \texttt{worker\_high} (mapped to a frontier reasoning
\texttt{gpt-5.5-2026-04-24-xhigh} configuration routed through a
local proxy at \texttt{localhost:4002}); the reviewer-sim sub-agent
uses \texttt{reviewer\_high} (mapped to a Claude-family
\texttt{claude-opus-4-7-thinking-high} configuration); the manager
uses \texttt{manager\_high} (also a Claude-family configuration but
distinct from the reviewer profile to avoid sharing prompt history).
The manager is woken by \texttt{launchd} every five minutes; the
reviewer-sim is invoked on-demand by the manager
(Algorithm~\ref{alg:tick}, line 16) when the worker's diff touches a
claim-bearing region. Each worker operates inside its own git
worktree (\texttt{paper/neurips2026-draft-worker-<id>}); commits
accumulate on the worker branch and never auto-merge to the trunk.

\paragraph{Honest disclosure of which loop ran when.}
We distinguish two operating regimes. The
\emph{rules-only} regime follows the per-paper hard rules
(\texttt{rules/paper\_c.md}, ten numbered rules HR-A--HR-J) by hand
and invokes the lit-review sub-agent on the seed bibliography; this
is the regime under which paper~C v1 was produced. The
\emph{full-loop} regime runs Algorithm~\ref{alg:tick} continuously,
with the manager auditing worker commits; paper~A and paper~B operate
in this regime on a separate compute host, with audit data not yet
available at submission time. The empirical claims in
\S\ref{sec:case-aggregate}--\S\ref{sec:case-citation} are computed on
the rules-only regime; eliding this distinction would itself be a
claim-drift violation under \S\ref{sec:discipline}.

\subsection{System Engineering Audit}
\label{sec:case-aggregate}

The complete engineering history of the discipline-mechanism
implementation is summarised in Table~\ref{tab:engineering}. The
fourteen commits spanning the development of this paper and the four
sub-agent modules add up to \(5{,}969\) lines of Python and Bash and
\(13/13\) passing hermetic self-tests. Two numbers from this table are worth highlighting.

\begin{table}[t]
  \caption{System engineering audit on the public release
    \texttt{main} branch (commit range: see Git history). Source:
    \texttt{git log --oneline} + \texttt{git log --shortstat}.}
  \label{tab:engineering}
  \centering
  \small
  \begin{tabular}{l r}
    \toprule
    \textbf{Quantity} & \textbf{Count} \\
    \midrule
    Total commits on branch                                   & 14 \\
    Cumulative line insertions across branch                  & $\sim 6{,}500$ \\
    Total source lines of system (Python + Bash)              &  5{,}969 \\
    Lines in the central CLI \texttt{mgr}                      &     770 \\
    Largest single sub-agent (\texttt{lib/lit\_review.py})    &     858 \\
    Smallest single sub-agent (\texttt{lib/lark\_notify.py})  &     157 \\
    Sub-agent contract document                                &     595 \\
    Hermetic self-tests passing (initial / final)             &  $8/8 \to 13/13$ \\
    Sub-agents shipped in parallel (single integration day)   &       4 \\
    \texttt{[BLOCKED-DECISION]} commits raised by humans      &       0 \\
    \texttt{[BLOCKED-DECISION]} commits raised by workers      &  pending\(^\dagger\) \\
    Bugs surfaced by self-test that would otherwise have shipped & 1 \\
    \bottomrule
  \end{tabular}

  \medskip
  \footnotesize
  \(^\dagger\)~Worker C (this paper) operates under the
  \texttt{[BLOCKED-DECISION]} discipline of \S\ref{sec:discipline-blocked}
  but at the time of writing the worker has not yet completed its first
  full draft cycle, so the count is reported as ``pending'' rather than
  ``zero'' to avoid the false-discipline reading.
\end{table}

\textbf{Sub-agent parallel ship.}
The sub-agent contract pattern (\S\ref{sec:arch-contract}) reduced the
implementation wall-clock for the four discipline modules from a
hand-estimated 3--4 work-days of sequential single-developer effort to
under thirty minutes of parallel LLM-driven implementation, a wall-clock
reduction of approximately two orders of magnitude. The
\(13/13\) extended smoke test (\(8/8 \to 13/13\) over the integration
of \texttt{scripts/\_invoke\_manager.py}) was the single integration
gate, and we observed zero post-integration regressions.

\textbf{One real bug auto-caught.}
While auditing the manager-backend dispatcher
(\texttt{scripts/\_invoke\_manager.py}, commit \texttt{c9ceeef}),
the contract self-test surfaced the following defect: when the
configured CLI executable did not exist on \texttt{PATH}, the
dispatcher raised an uncaught \texttt{FileNotFoundError} traceback to
stderr instead of returning the contracted single-line JSON to stdout,
which would have caused \texttt{tick.sh} to write a malformed
escalation entry on every failed tick. The fix (catch the exception,
write a friendly note to the captured stdout file, return exit code
\(127\), emit the contracted JSON) was made in the same commit and
is regression-tested by four explicit error-path probes whose output
is verified against the dispatcher's contract. The bug is exactly the
class of silent-failure mode that \S\ref{sec:discipline-sentinel}
designs the Sentinel watchdog to catch in deployment; finding it
pre-deployment via the hermetic self-test is consistent with the
discipline value of the contract pattern.

\subsection{Cross-Model Adversarial Ablation}
\label{sec:case-ablation}

We pre-registered the cross-model ablation
of claim~C3 with three conditions: (i) same model, same backend
(self-review baseline); (ii) different model, same backend
(model-only ablation); (iii) full cross-API-model configuration in
which worker and reviewer-sim run on \emph{different} codex profiles
mapped to different underlying API model families
(e.g.\ \texttt{worker\_high} $\to$ gpt-5.5 vs.\ \texttt{reviewer\_high}
$\to$ claude-opus). The protocol is to inject an Abstract
claim-softening commit (e.g., changing ``the gap is not a token-budget
problem and not a model-capability problem'' to ``the gap may be partly
a token-budget problem'') and to count the fraction of injections in
which the manager LLM cancels the commit and escalates to the human
within two ticks.

At the time of paper submission, the ablation runs against conditions
(i) and (iii) are scheduled but have not completed in the audit log;
to remain compliant with HR-F we report this as
\texttt{\textbackslash cite\{NUMBERS-PENDING-AUDIT-FETCH\}} rather than
imputing a number, and we will publish the completed table to the
open-source repository before the camera-ready deadline. We note the
qualitative observation, however, that across the development of this
paper itself (in which the manager role was played by a Claude-family
codex profile and the worker role was played by a separate Claude-family
codex profile, with the GPT-5.5 reviewer-sim profile not yet wired
in), the manager flagged two distinct claim-softening drafts proposed
by the worker that were subsequently revised; the same drafts evaluated
under condition (i) by hand-replay produced no flag. We treat this as
sub-quantitative evidence consistent with C3 and defer the
quantitative table to the camera-ready.

\subsection{Lit-Review Citation Verification}
\label{sec:case-citation}

The lit-review sub-agent was invoked once on the seed
\texttt{paper\_meta/references.bib} of this paper prior to substantive
drafting and a second time after a reviewer-anticipation pass
identified a missing entry. Of the nine bibliography entries
ultimately included, all nine had their primary metadata fields
(title, authors, year, arXiv id or canonical URL) verified against
arXiv, ACL Anthology, GitHub, or the publishing organization's blog;
three had defects (one inserted by the seed bibliography drafter,
one minor canonical-name issue, and one missing-entry pattern that
required a second pass).
Table~\ref{tab:litreview} summarises the verification outcomes.

\begin{table}[t]
  \caption{Outcome of \texttt{mgr lit-review paper\_c} on the seed
    bibliography of this paper. ``Defect found'' means the seed entry
    had a non-trivial metadata error; the dispatcher raised a
    \texttt{[BLOCKED-DECISION]}-equivalent advisory rather than
    silently rewriting the bibliography.}
  \label{tab:litreview}
  \centering
  \footnotesize
  \begin{tabular}{l c c p{0.38\textwidth}}
    \toprule
    \textbf{Citation key} & \textbf{Verified} & \textbf{Defect} & \textbf{Defect detail (if any)} \\
    \midrule
    \texttt{song2026paperorchestra}        & yes & no  & --- \\
    \texttt{sibyl2026system}               & yes & yes & seeded with fictitious ``renamed from FARS'' note; corrected to acknowledge separate provenance \\
    \texttt{fars2026analemma}              & yes & yes & added in second pass after the FARS-vs-us positioning gap was identified \\
    \texttt{yamada2025aiscientistv2}       & yes & no  & --- \\
    \texttt{karpathy2026autoresearch}      & yes & yes & canonical repo name is lowercase \texttt{autoresearch}, not \texttt{AutoResearch} \\
    \texttt{schmidgall2025agentlaboratory} & yes & no  & --- \\
    \texttt{lu2024aiscientist}             & yes & no  & --- \\
    \texttt{tang2025airesearcher}          & yes & no  & --- \\
    \texttt{miyai2026paperreconstruction}  & yes & no  & --- \\
    \bottomrule
  \end{tabular}
\end{table}

The three surfaced defects are minor in isolation but illustrative
of the failure modes the lit-review sub-agent is targeted at: a
fictitious provenance claim
(\texttt{sibyl2026system} seeded with a ``renamed from FARS'' note
that the verification pass refuted by reading the project README and
the FARS primary source); a wrong canonical name
(\texttt{karpathy/AutoResearch} versus the actual
\texttt{karpathy/autoresearch}); and a missing-entry pattern
(\texttt{fars2026analemma} was absent from the seed despite being
the most quantitatively important comparison point in
\S\ref{sec:related}, surfaced when a reviewer-anticipation pass
asked which scale-throughput baseline the paper engaged). All three
would have produced reviewer-visible defects in the camera-ready
had they shipped uncorrected.

The Sentinel watchdog of \S\ref{sec:discipline-sentinel} was
specified in response to a real incident on the deployment host
(\texttt{gpu\_develop}, 2026-05-02) in which the LLM CLI in use at
the time produced an identical exit-code-1 failure on every
five-minute tick for approximately twenty-four hours before the
deployment owner manually noticed the pattern in
\texttt{state/escalations.jsonl}. (The CLI in question was
cursor-agent under an early prototype; the codex-only loop has the
same failure mode and the same need for a watchdog.)
Sentinel's same-cause-streak heuristic
(\S\ref{sec:discipline-sentinel}, item~ii) is parameterised to fire
after four consecutive identical failures; the same incident under
the watchdog would have produced a Lark notification within twenty
minutes (a reduction by a factor of approximately seventy). The full
incident reconstruction and rendered Lark message are in
Appendix~\ref{app:sentinel-msg}.
